%% toc: P15 161 161.1,161.2,161.4,161.5
% Almost-sure, probability, distribution and mean-power convergence compared (Durrett, secs. 2.2, 2.3, 3.2, 4.4, 4.6).
\chapter{Almost-Sure, Probability, Distribution and Mean-Power Convergence}
\label{sec:p15-w02b-p15-013-modes}

Four ways in which a sequence of random variables can approach a limit
recur through this part: convergence almost surely, in probability, in
distribution, and in \(p\)th mean. This chapter compares them as a
whole. It states which implications hold, under which side conditions
and for which exponents \(p\), proves each of them, and gives for every
implication that fails a model in which it fails. The mean-power
exponent is allowed to be any real \(p>0\). The three ranges
\(0<p<1\), \(p=1\) and \(1<p<\infty\) behave differently in their
metric structure, but the comparison of modes turns out to be the same
in all three, and each theorem below says so explicitly.

The source is Durrett's author version of \emph{Probability: Theory and
Examples}, dated 11 January 2019. Durrett defines convergence in
probability at the start of Section~2.2 \cite[p.~56]{DurrettPTEAuthor2019}
and proves, for every \(p>0\), that \(\mathbb E|Z_n|^p\to0\) forces
\(Z_n\to0\) in probability \cite[Lemma~2.2.2, p.~57]{DurrettPTEAuthor2019}.
His Theorems~2.3.2--2.3.4 relate convergence in probability to
almost-sure convergence along subsequences, and the remark after them
draws the consequence that almost-sure convergence comes from no
topology \cite[pp.~68--69]{DurrettPTEAuthor2019}. Exercise~2.3.11 is his
model separating the two \cite[p.~75]{DurrettPTEAuthor2019}; Exercise~3.2.12
links probability and distribution
\cite[p.~124]{DurrettPTEAuthor2019}; and Theorems~4.6.2--4.6.3 give the
uniform-integrability criterion for convergence in mean
\cite[pp.~245--246]{DurrettPTEAuthor2019}. Durrett states that criterion
only for the first mean. Its extension to every \(p>0\), the inclusion
of \(L^p\) in \(L^q\) for \(q<p\) in the range below one, and the
assembled list of countermodels are supplied here.

\emph{Inputs.} A probability space \((\Omega,\mathcal F,\mathbb P)\);
finite real random variables \(X_n\) (\(n\geq1\)) and \(X\) on it; and,
for the mean-power statements, an exponent \(0<p<\infty\). Convergence in
distribution alone needs no common space; it is compared with the
other modes only for variables on one space.
\emph{Conclusions.} The implication diagram of
\cref{thm:p15-w02b-p15-013-comparison}: almost-sure convergence and
\(p\)th-mean convergence each imply convergence in probability, which
implies convergence in distribution; \(p\)th-mean convergence implies
\(q\)th-mean convergence for \(0<q<p\); convergence in probability
together with uniform integrability of \(|X_n|^p\) is equivalent to
\(p\)th-mean convergence (\cref{thm:p15-w02b-p15-013-vitali}); and no
further implication holds without additional hypotheses
(\cref{thm:p15-w02b-p15-013-countermodels}).

\section{The four modes and the mean-power ranges}
\label{sec:p15-w02b-p15-013-definitions}

\begin{hbdefinition}{def:p15-w02b-p15-013-modes}{Four modes of convergence}
Let \(X_n\) (\(n\geq1\)) and \(X\) be finite real random variables on
\((\Omega,\mathcal F,\mathbb P)\).
\begin{enumerate}
\item \(X_n\to X\) \emph{almost surely} (a.s.) if the event
      \[
       \{\omega:X_n(\omega)\to X(\omega)\}=
       \bigcap_{r\geq1}\bigcup_{N\geq1}\bigcap_{n\geq N}
       \{|X_n-X|\leq1/r\}
      \]
      has probability one. The displayed countable combination shows
      that it is an event; compare \cref{thm:p15-event-limit-calculus}.
\item \(X_n\to X\) \emph{in probability} if
      \(\mathbb P(|X_n-X|>\varepsilon)\to0\) for every
      \(\varepsilon>0\). This is Durrett's definition
      \cite[p.~56]{DurrettPTEAuthor2019} and the probability case of
      convergence in measure in \cref{def:p15-lp-convergence}.
\item \(X_n\Rightarrow X\), \emph{in distribution}, if the CDFs
      converge at every continuity point of the limiting CDF, as in
      \cref{def:p15-weak-law-convergence}. Only the laws matter.
\item For \(0<p<\infty\), put
      \(\|X\|_p\defeq(\mathbb E|X|^p)^{1/p}\in[0,\infty]\). The set
      \(L^p(\mathbb P)\) consists of the classes, under a.s. equality,
      of finite real random variables with \(\mathbb E|X|^p<\infty\).
      \(X_n\to X\) \emph{in \(p\)th mean}, or \emph{in \(L^p\)}, if
      \(\mathbb E|X_n-X|^p\to0\).
\end{enumerate}
For \(p=1\) the fourth item is the \(L^1\) convergence of
\cref{def:p15-lp-convergence}; for \(1\leq p<\infty\) its size is the
norm of \cref{def:p15-lp-conditional-spaces} and
\cref{def:p15-w02b-p15-001-lp-space}. For \(0<p<1\) the same formula
is used, although it is not a norm (\cref{thm:p15-w02b-p15-013-power}).
The definition of \(p\)th-mean convergence does not require
\(X_n\) or \(X\) to lie in \(L^p\); the theorem just cited shows that
once \(\mathbb E|X_n-X|^p\) is finite, \(X_n\) lies in \(L^p\) exactly
when \(X\) does.
\index{convergence!almost sure}\index{convergence!in probability}
\index{convergence!in pth mean@in \(p\)th mean}
\end{hbdefinition}

\begin{hbconvention}{def:p15-w02b-p15-013-ranges}{The three mean-power ranges}
The range \(0<p<1\) is the range in which \(t\mapsto t^p\) is concave on
\([0,\infty)\); there \(|a+b|^p\leq|a|^p+|b|^p\), the quantity
\(\mathbb E|X-Y|^p\) itself is a metric, and \(\|\cdot\|_p\) fails the
triangle inequality. At \(p=1\) the size is the \(L^1\) norm and the
expectation is defined on the whole space. For \(1<p<\infty\),
\(t\mapsto t^p\) is convex, \(\|\cdot\|_p\) is a norm by Minkowski's
inequality, and H\"older's inequality pairs \(L^p\) with
\(L^{p/(p-1)}\). The constant
\[
 c_p\defeq\max(1,2^{p-1})
\]
records the difference in the pointwise inequality for sums.
\end{hbconvention}

\begin{hbtheorem}{thm:p15-w02b-p15-013-power}{proved=annex:p15-w02b-p15-013-power}{Power inequalities and the spaces \(L^p\) for every \(p>0\)}
\par\nobreak
Let \(a,b\) be real and \(0<p<\infty\).
\begin{enumerate}
\item If \(0<p\leq1\), then \(|a+b|^p\leq|a|^p+|b|^p\) and
      \(\bigl||a|^p-|b|^p\bigr|\leq|a-b|^p\).
\item If \(1\leq p<\infty\), then
      \(|a+b|^p\leq2^{p-1}(|a|^p+|b|^p)\).
\item Hence \(|a+b|^p\leq c_p(|a|^p+|b|^p)\) for every \(p>0\), and
      \(L^p(\mathbb P)\) is a real vector space under pointwise
      operations.
\item For \(0<p<1\), \(d_p(X,Y)\defeq\mathbb E|X-Y|^p\) is a metric on
      \(L^p(\mathbb P)\). For \(1\leq p<\infty\),
      \(\|X-Y\|_p\) is a metric on it. In each range, \(X_n\to X\) in
      \(L^p\) means convergence in the corresponding metric whenever all
      variables lie in \(L^p\).
\item If \(X_n\to X\) and \(X_n\to Y\) in \(L^p\), then \(X=Y\) a.s.
\end{enumerate}
\end{hbtheorem}

\begin{hbtheorem}{thm:p15-w02b-p15-013-monotone}{proved=annex:p15-w02b-p15-013-monotone}{Smaller exponents: \(L^p\) inside \(L^q\)}
For \(0<q<p<\infty\) and every finite real random variable \(X\),
\[
 \|X\|_q\leq\|X\|_p\qquad\text{in }[0,\infty].
\]
Consequently \(L^p(\mathbb P)\subseteq L^q(\mathbb P)\), and
\(X_n\to X\) in \(L^p\) implies \(X_n\to X\) in \(L^q\). In particular,
for \(p\geq1\), convergence in \(L^p\) implies convergence in \(L^1\);
if moreover \(X\in L^1\), then \(X_n\in L^1\) for all large \(n\) and
\(\mathbb EX_n\to\mathbb EX\). The pair \(q,p\) may lie in
the same range or in different ranges.
\end{hbtheorem}

The inequality is proved from H\"older's inequality applied to
\(|X|^q\) and the constant one, with exponent \(p/q>1\)
(\cref{thm:p15-holder-integral}; \cite[Theorem~1.6.3, p.~29]{DurrettPTEAuthor2019}).
It needs \(\mathbb P(\Omega)=1\): \cref{thm:p15-w02b-p15-013-finite-measure}
shows that it fails for an infinite measure.

\section{Almost-sure convergence and subsequences}
\label{sec:p15-w02b-p15-013-as}

\begin{hbtheorem}{thm:p15-w02b-p15-013-as-probability}{proved=annex:p15-w02b-p15-013-as-probability}{Almost-sure convergence, subsequences and the absence of a topology}
\par\nobreak
Let \(X_n\), \(X\), \(Y\) be finite real random variables.
\begin{enumerate}
\item If \(X_n\to X\) a.s., then \(X_n\to X\) in probability.
\item If \(X_n\to X\) and \(X_n\to Y\) in probability, then \(X=Y\)
      a.s.
\item \(X_n\to X\) in probability if and only if every subsequence
      \((X_{n(m)})\) has a further subsequence \((X_{n(m_k)})\)
      converging to \(X\) a.s.
\item If \(X_n\to X\) in probability and
      \(g:\mathbb R\to\mathbb R\) is continuous, then
      \(g(X_n)\to g(X)\) in probability.
\item In a topological space, if every subsequence of a sequence
      \((y_n)\) has a further subsequence converging to \(y\), then
      \(y_n\to y\).
\item There is no topology on the set of finite real random variables
      on the uniform probability space of
      \cref{thm:p15-lebesgue-unit-interval}, or on the set of their
      a.s.\ classes, whose convergent sequences are exactly the
      almost surely convergent ones, with the same limits.
\item If \(\Omega\) is countable and \(\mathcal F\) consists of all
      its subsets, then \(X_n\to X\) in probability implies
      \(X_n\to X\) a.s.
\end{enumerate}
\end{hbtheorem}

Items~(3)--(5) are Durrett's Theorems~2.3.2, 2.3.4 and~2.3.3, and
item~(6) is the remark that follows them
\cite[pp.~68--69]{DurrettPTEAuthor2019}. Item~(1) is the special case of
Durrett's Theorem~2.3.2 in which the whole sequence converges a.s.;
the proof here is direct. One direction of item~(3) is
\cref{thm:p15-probability-l1-tools}~(2), which is reused. Item~(6)
is proved with the typewriter sequence of
\cref{thm:p15-w02b-p15-013-countermodels} in place of Durrett's
pointer to his exercises. Item~(7) is his Exercise~2.3.12
\cite[p.~75]{DurrettPTEAuthor2019}: on a countable space every point of
positive probability must eventually see the limit, so the typewriter
phenomenon of model~(2) needs an uncountable sample space with
null singletons.

\section{The named comparison}
\label{sec:p15-w02b-p15-013-comparison}

\begin{hbtheorem}{thm:p15-w02b-p15-013-comparison}{proved=annex:p15-w02b-p15-013-comparison}{The comparison of the four modes in every mean-power range}
Let \(X_n,X\) be finite real random variables on one probability space
and let \(0<q<p<\infty\). The following hold in each of the ranges
\(0<p<1\), \(p=1\) and \(1<p<\infty\), and for \(q\) in any range
below \(p\).
\begin{enumerate}
\item \emph{Almost sure to probability.} \(X_n\to X\) a.s. implies
      \(X_n\to X\) in probability.
\item \emph{Mean to probability.} \(X_n\to X\) in \(L^p\) implies
      \(X_n\to X\) in probability, with
      \(\mathbb P(|X_n-X|\geq\varepsilon)\leq
      \varepsilon^{-p}\mathbb E|X_n-X|^p\) for \(\varepsilon>0\)
      \cite[Lemma~2.2.2, p.~57]{DurrettPTEAuthor2019}.
\item \emph{Higher to lower mean.} \(X_n\to X\) in \(L^p\) implies
      \(X_n\to X\) in \(L^q\).
\item \emph{Probability to distribution.} \(X_n\to X\) in probability
      implies \(X_n\Rightarrow X\). Conversely, if \(X_n\Rightarrow c\)
      for a real constant \(c\), then \(X_n\to c\) in probability
      \cite[Exercise~3.2.12, p.~124]{DurrettPTEAuthor2019}.
\item \emph{Probability to almost sure, partially.} If \(X_n\to X\) in
      probability, a subsequence converges to \(X\) a.s.
\item \emph{Probability to mean.} If \(X_n\to X\) in probability, then
      \(X_n\to X\) in \(L^p\) if and only if \(X_n\in L^p\) for every
      \(n\) and \((|X_n|^p)_{n\geq1}\) is uniformly integrable
      (\cref{thm:p15-w02b-p15-013-vitali}). Two sufficient conditions
      are domination, \(|X_n|\leq Z\) a.s.\ with \(Z\in L^p\), and a
      bound \(\sup_n\mathbb E|X_n|^r<\infty\) for some \(r>p\).
\end{enumerate}
The implications (1)--(6), their compositions, and the trivial ones
from a mode to itself are the only implications among the modes
almost surely, in probability, in distribution, in \(L^q\) and in
\(L^p\) that hold for all sequences: every other implication between
two of these modes fails in a model of
\cref{thm:p15-w02b-p15-013-countermodels}, for each \(p\) and \(q\).
\end{hbtheorem}

\begin{hbtheorem}{thm:p15-w02b-p15-013-vitali}{proved=annex:p15-w02b-p15-013-vitali}{Uniform integrability and \(L^p\) convergence}
Let \(0<p<\infty\), let \(X_n\in L^p(\mathbb P)\) for every \(n\),
and let \(X_n\to X\) in probability, with \(X\) finite, real and
measurable. The following are equivalent:
\begin{enumerate}
\item \((|X_n|^p)_{n\geq1}\) is uniformly integrable;
\item \(X\in L^p(\mathbb P)\) and \(\mathbb E|X_n-X|^p\to0\);
\item \(\mathbb E|X_n|^p\to\mathbb E|X|^p<\infty\).
\end{enumerate}
In particular, if \(|X_n|\leq Z\) a.s.\ for every \(n\) with
\(Z\in L^p\), or if \(\sup_n\mathbb E|X_n|^r<\infty\) for some
\(r>p\), then \(X_n\to X\) in \(L^p\).
For \(p=1\) this is \cref{thm:p15-vitali-convergence}, which is
Durrett's Theorem~4.6.3 \cite[pp.~245--246]{DurrettPTEAuthor2019}.
\end{hbtheorem}

For \(p>1\) and martingales, \cref{thm:p15-martingale-lp-convergence}
is a special case of the sufficient condition
\(\sup_n\mathbb E|X_n|^r<\infty\) only after the almost-sure limit has
been found; there the martingale structure supplies that limit, which
the present theorem takes as a hypothesis
\cite[Theorem~4.4.6, p.~237]{DurrettPTEAuthor2019}. The superlinear
moment criterion used for the second sufficient condition is
\cref{thm:p15-ui-moment-criteria}, Durrett's Theorem~4.6.2
\cite[p.~245]{DurrettPTEAuthor2019}, applied to \(|X_n|^p\) with
\(\Phi(t)=t^{r/p}\).

\section{Countermodels for every failing converse}
\label{sec:p15-w02b-p15-013-countermodels}

All models except the third live on the uniform probability space of
\cref{thm:p15-lebesgue-unit-interval}, with sample space \((0,1)\) and
probability \(\mathbb U\).

\begin{hbtheorem}{thm:p15-w02b-p15-013-countermodels}{proved=annex:p15-w02b-p15-013-countermodels}{Separating models in each mean-power range}
\par\nobreak
Fix \(0<p<\infty\) and \(0<q<p\).
\begin{enumerate}
\item \emph{A spike.} Let \(Z^{(p)}_n=n^{1/p}\mathbf1_{(0,1/n)}\).
      Then \(Z^{(p)}_n\to0\) at every point of \((0,1)\), hence a.s.,
      in probability and in distribution; \(\mathbb E|Z^{(p)}_n|^q=
      n^{q/p-1}\to0\), so \(Z^{(p)}_n\to0\) in \(L^q\); but
      \(\mathbb E|Z^{(p)}_n|^p=1\) for every \(n\), and
      \((Z^{(p)}_n)\) converges in \(L^p\) to no random variable.
      The family \((|Z^{(p)}_n|^p)\) is bounded in \(L^1\) and is not
      uniformly integrable.
\item \emph{A typewriter.} For \(n=2^k+j\) with integers \(k\geq0\) and
      \(0\leq j<2^k\), let
      \(T_n=\mathbf1_{[j2^{-k},(j+1)2^{-k})\cap(0,1)}\). Then
      \(\mathbb E|T_n|^p=2^{-k}\to0\) for every \(p>0\), so
      \(T_n\to0\) in every \(L^p\), in probability and in distribution;
      but at every point of \((0,1)\) the sequence \(T_n\) takes both
      values zero and one infinitely often, so it converges at no point,
      and there is no random variable to which it converges a.s.
\item \emph{Independent rare indicators.} This is Durrett's
      Exercise~2.3.11 \cite[p.~75]{DurrettPTEAuthor2019}. Let
      \(B_1,B_2,\ldots\) be independent with \(\mathbb P(B_n=1)=p_n\) and
      \(\mathbb P(B_n=0)=1-p_n\), where \(0\leq p_n\leq1\). Then
      (i) \(B_n\to0\) in probability if and only if \(p_n\to0\), and
      this holds if and only if \(B_n\to0\) in \(L^p\), for any
      \(p>0\); (ii) \(B_n\to0\) a.s.\ if and only if
      \(\sum_np_n<\infty\). With \(p_n=1/n\), the sequence converges to
      zero in every \(L^p\) and in probability, and has no a.s.\ limit.
\item \emph{Sign reversal.} Let \(\zeta=2\mathbf1_{(0,1/2)}-1\),
      a fair sign, and \(X_n=-\zeta\) for every \(n\). Then
      \(X_n\Rightarrow\zeta\), but
      \(\mathbb P(|X_n-\zeta|>1)=1\) and
      \(\mathbb E|X_n-\zeta|^p=2^p\) for every \(n\). So
      \(X_n\) converges to \(\zeta\) neither in probability, nor
      in \(L^p\), nor a.s.
\item \emph{Below the first mean.} Let \(Y=2^k\) on
      \([2^{-k},2^{1-k})\), \(k\geq1\). For \(0<p<1\),
      \(\mathbb EY^p=1/(2^{1-p}-1)<\infty\), while \(\mathbb EY=\infty\).
      The variables \(Y/n\) converge to zero a.s.\ and in \(L^p\) for
      every \(0<p<1\), but none of them lies in \(L^1\), so
      \(\mathbb E|Y/n-0|=\infty\) for every \(n\).
\end{enumerate}
The failing implications are distributed as follows. Model~(1)
refutes almost sure \(\Rightarrow L^p\), probability \(\Rightarrow L^p\),
distribution \(\Rightarrow L^p\) and \(L^q\Rightarrow L^p\), and shows that a
bound on \(\mathbb E|X_n|^p\) alone does not give \(L^p\) convergence;
the same model with exponent \(q\), namely \(Z^{(q)}_n\), refutes
almost sure \(\Rightarrow L^q\) and probability \(\Rightarrow L^q\).
Models~(2) and~(3) refute \(L^p\Rightarrow\) almost sure,
probability \(\Rightarrow\) almost sure and distribution \(\Rightarrow\) almost
sure. Model~(4) refutes distribution \(\Rightarrow\) probability and,
again, distribution \(\Rightarrow L^p\) and distribution \(\Rightarrow\) almost
sure, even with all variables bounded on one space. Model~(5) shows,
inside the range \(0<p<1\), that convergence in every \(L^p\) with
\(p<1\) does not even make the variables integrable.
\end{hbtheorem}

At \(p=1\), model~(1) is \(Z^{(1)}_n=n\mathbf1_{(0,1/n)}\), the source's
example of \cref{thm:p15-ui-counterexamples}~(1) and the concentrating
functions of \cref{def:p15-convergence-failures}; at \(p=2\) it is
\(\sqrt n\,\mathbf1_{(0,1/n)}\); at \(p=1/2\) it is
\(n^2\mathbf1_{(0,1/n)}\). For submartingales and \(p>1\) the same
phenomenon appears in \cref{thm:p15-maximal-convergence-boundaries}~(1),
and \cref{thm:p15-weak-convergence-boundaries}~(3) records the
probability part of model~(4).

\begin{hbtheorem}{thm:p15-w02b-p15-013-finite-measure}{proved=annex:p15-w02b-p15-013-finite-measure}{Where the total mass one is used}
On \(\mathbb R\) with Lebesgue measure \(\lambda\):
\begin{enumerate}
\item \(f_n=\mathbf1_{[n,\infty)}\) converges to zero at every point,
      but \(\lambda(\{|f_n|>1/2\})=\infty\) for every \(n\); so
      convergence everywhere does not imply convergence in measure.
\item \(f=\sum_{k\geq1}k^{-1}\mathbf1_{[k,k+1)}\) satisfies
      \(\int f^2\,d\lambda<\infty\) and \(\int f\,d\lambda=\infty\);
      so \(L^2(\lambda)\not\subseteq L^1(\lambda)\), and the
      inequality of \cref{thm:p15-w02b-p15-013-monotone} fails.
\end{enumerate}
\end{hbtheorem}

\section{Procedure, worked example and limits}
\label{sec:p15-w02b-p15-013-use}

\begin{hbmethod}{def:p15-w02b-p15-013-method}{Determining the modes in which a sequence converges}
Given finite real \(X_n\) on one probability space and a candidate
limit \(X\):
\begin{enumerate}
\item Settle convergence in probability first. Use a direct tail
      estimate, an a.s.\ argument (\cref{thm:p15-w02b-p15-013-as-probability}~(1)),
      a moment bound \(\mathbb E|X_n-X|^p\to0\) for some \(p>0\)
      (\cref{thm:p15-w02b-p15-013-comparison}~(2)), or, for a constant
      \(X\), weak convergence (item~(4)). If it fails, no mode except
      possibly convergence in distribution can hold, and a candidate
      limit for any of them is unique a.s.
\item For each exponent \(p\) of interest, test uniform integrability
      of \(|X_n|^p\) with \cref{thm:p15-w02b-p15-013-vitali}: a
      dominating \(Z\in L^p\), a bound on a higher moment, or the
      equality \(\mathbb E|X_n|^p\to\mathbb E|X|^p<\infty\). The side
      conditions are \(X_n\in L^p\) for each \(n\) and convergence in
      probability; uniform integrability for one \(p\) gives it for every
      smaller exponent, by \cref{thm:p15-w02b-p15-013-monotone} applied
      to the conclusion.
\item For almost-sure convergence, convergence in probability supplies
      only a subsequence. Use a summable tail bound and
      \cref{thm:p15-borel-cantelli-first}, a monotonicity or martingale
      argument, or exhibit failure with the second Borel--Cantelli lemma
      (\cref{thm:p15-borel-cantelli-second}) as in model~(3).
\item Convergence in distribution transfers to the other modes only
      toward a constant, and only through convergence in probability.
\end{enumerate}
\end{hbmethod}

\begin{hbexample}{def:p15-w02b-p15-013-worked}{A spike that converges only below one half}
On the uniform space let \(X_n=n^2\mathbf1_{(0,1/n)}\). At every
\(u\in(0,1)\), \(X_n(u)=0\) once \(n\geq1/u\), so \(X_n\to0\)
everywhere, hence a.s., in probability and in distribution. For
\(r>0\),
\[
 \mathbb E|X_n|^r=n^{2r}\cdot\frac1n=n^{2r-1},
\]
which tends to zero exactly when \(r<1/2\), equals one at \(r=1/2\),
and tends to infinity for \(r>1/2\). Thus \(X_n\to0\) in \(L^r\) for
\(0<r<1/2\) and in no \(L^r\) with \(r\geq1/2\); in particular not in
\(L^1\) or in \(L^2\). The procedure finds the same answer without the
explicit limit. For \(r=1/4\), \(\mathbb E(|X_n|^{1/4})^2=
\mathbb E|X_n|^{1/2}=1\), so the family \((|X_n|^{1/4})\) is bounded in
\(L^2\) and therefore uniformly integrable by
\cref{thm:p15-ui-moment-criteria}; \cref{thm:p15-w02b-p15-013-vitali}
then gives \(L^{1/4}\) convergence. For \(r=1/2\), the variables
\(|X_n|^{1/2}=n\mathbf1_{(0,1/n)}\) are the non-uniformly-integrable
family of \cref{thm:p15-ui-counterexamples}~(1), and the same theorem
rules out \(L^{1/2}\) convergence; it rules out every \(r\geq1/2\) by
\cref{thm:p15-w02b-p15-013-monotone}. This is model~(1) of
\cref{thm:p15-w02b-p15-013-countermodels} with \(p=1/2\).
\end{hbexample}

\begin{hbwarning}{def:p15-w02b-p15-013-limits}{What the comparison does not say}
\par\nobreak
\begin{enumerate}
\item The implication from almost-sure convergence to convergence in
      probability, and the inclusion of \(L^p\) in \(L^q\) for
      \(q<p\), use total mass one; both fail for Lebesgue measure on
      \(\mathbb R\) (\cref{thm:p15-w02b-p15-013-finite-measure}).
\item For \(0<p<1\), \(\|\cdot\|_p\) is not a norm
      (\cref{def:p15-w02b-p15-001-lp-limits}~(2)); the metric is
      \(\mathbb E|X-Y|^p\) without the root. Expectation need not be
      defined on \(L^p\) in this range (model~(5)), so
      \(L^p\) convergence for \(p<1\) never by itself yields convergence
      of means.
\item Bounded \(p\)th moments with convergence in probability do not
      give \(L^p\) convergence (model~(1)); a bound on a strictly higher
      moment does. Equality of limits in distribution says nothing about
      the joint behavior of \(X_n\) and \(X\) (model~(4)).
\item Convergence in \(L^p\) for all \(p<\infty\) together with a
      uniform bound \(|X_n|\leq1\) still does not give a.s.\
      convergence (model~(2)).
\end{enumerate}
\end{hbwarning}

\begin{proofboundary}{def:p15-w02b-p15-013-relations}{Neighboring chapters}
Weak convergence itself, its metrics, tightness and the continuous
mapping theorem are developed in \cref{sec:p15-weak-convergence-of-laws},
where \cref{thm:p15-weak-metric-premises} proves the probability bridge
used in \cref{thm:p15-w02b-p15-013-comparison}~(4). Uniform
integrability, the case \(p=1\) of \cref{thm:p15-w02b-p15-013-vitali},
and its martingale consequences are in
\cref{sec:p15-uniform-integrability-and-conditional-limits}. The
\(L^p\) norm laws for \(p\geq1\) are \cref{thm:p15-w02b-p15-001-lp-results}
and \cref{thm:p15-conditional-lp}. Laws of large numbers supply the
standing examples of convergence in probability
(\cref{thm:p15-direct-variance-weak-laws}) and almost surely
(\cref{thm:p15-etemadi-strong-law}).
\end{proofboundary}
